\section{Skill Challenges}
Skill Challenges are a system for non-combat encounters and challenges that the PCs undertake, they are especially useful when the situation cannot be resolved by just a single roll.
Skill Challenges can be used in many settings, like for travel, infiltrations, long social encounters, and so on.
The general approach of a Skill Challenge is that all involved PCs will make a roll based on how they are contributing to the goal and based on the result of all the rolls the players and GM gain Advantage that can be used to mark \emph{progress} on tracks, which may reflect either progress towards a goal or getting closer to some complication.
A \emph{track} represents some goal or complication that may occur and has some length established by the GM.
The more times a track has been marked the closer it gets to completion, when it reaches its length it is finished: the goal is completed or the complication now comes into play.
Some tracks can also be \emph{progressive}, meaning that as the track moves further along, it begins to impact the narrative.
For example, during travel a progressive track is used to track the party's progress and as they get further along in the track it may change details like the environment they are in or the weather.
Similarly, during an infiltration there might be a progressive complication track representing being caught and as that track ticks up the presence of guards may increase as they become suspicious, which might in turn make future rolls to avoid detection harder.
Additionally, some tracks may be \emph{linked}, meaning that when one track is ticked the linked tracks are automatically ticked as well which can represent how progress and consequences can be tied together.
For example, during travel the exhaustion track is linked to the progress track so that when the party makes more progress they get closer to exhaustion as well.
Finally, in addition to these tracks, a Skill Challenge generally has some kind of \emph{cost} which could be a literal cost in the game or some event that will occur each round of rolls, which reflects the time the in-game events are taking and how they may impact the scenario.
As an example, while traveling the GM may decide that each round represents a day of travel and that each day there will be an encounter roll to determine what, if anything, the party encounters.
Alternatively, if the party is working on some goal against an adversary the cost each round may represent things that adversary does to make the party's work more difficult.
Similarly, each round there may be some \emph{maintenance} that occurs which may update the state of certain tracks, such as always increasing a track representing something that always occurs or resetting some track that represents goals or complications for an individual round.

\subsection{Setting Up a Skill Challenge}
When a GM determines that a skill challenge is appropriate they will determine the following for the challenge; they may consult with the players in setting up the challenge.
\begin{enumerate}
  \item Determine what the PCs' goals are, if there are multiple distinct goals these may each be assigned a distinct track but if all the goals are related they may be assigned to a single track.
    Based on the difficulty of the goal(s) of each track, determine the length of each track.
  \item Determine any consequences that may occur, assign them tracks, and determine the lengths of those tracks.
  \item Determine other details of the tracks, in particular:
  \begin{enumerate}[a.]
    \item Determine which (if any) tracks are \emph{progressive} and describe how progress on that track will impact the narrative or challenge.
    \item Determine which (if any) tracks are \emph{linked} and the direction (i.e., ticking which track automatically ticks the other track).
  \end{enumerate}
  \item Determine any \emph{maintenance} that occur each round, most commonly this might include:
  \begin{itemize}
    \item Resetting a track (like a scouting tracking) each round.
    \item Automatically ticking a track (like for progress in travel) each round.
  \end{itemize}
  \item Determine any \emph{cost} that occurs each round and the consequences, such as encounters (or encounter rolls).
\end{enumerate}

\subsection{Running a Skill Challenge}
To run a skill challenge, the group follows the following steps:
\begin{enumerate}
  \item Any \emph{maintenance} that the GM has determine occurs will occur at the beginning of the round, potentially updating tracks.
  \item The GM then asks each player what their character is doing to help with the challenge and has them make an appropriate roll to help.
    Like a Cooperative Roll the number of successes and failures are totaled, and based on the result the GM and/or players will gain a number of Advantage that can be saved or spent to advance any of the tracks of the challenge.
  \begin{itemize}
    \item If the total of the successes and failures is $-2$ or less, the GM gains Advantage equal to half of absolute value of the total (rounded towards zero).
    \item If the total is between $-1$ and $+1$, the players and GM both gain 1 Advantage.
    \item If the total is $+2$ or greater, the players gain Advantage equal to half of the total (rounded towards zero).
  \end{itemize}
  \item The GM spend Advantage (either those they gained from the roll or any others they have) and the players spend Advantage (only those they gained from the roll) to advance any tracks they want, or to push back any tracks they wish (for instance to avoid a complication).
    Then, the outcomes of any progressive tracks or tracks that have been completed are resolved (potentially ending the challenge if the players decide they have completed their goals sufficiently).
  \item Finally, any \emph{cost} is resolved at the end of the round and once resolved the next round of the challenge commences.
\end{enumerate}

\subsection{Example: Travel}
As an example of how to setup a skill challenge, this section describes a general form for skill challenges representing a party's travels.
\begin{enumerate}
  \item The main goal track is reaching the PCs' destination, each tick on the track should represent an easy day's travel and the party may be expected to complete two or three (or even more) ticks a day.
    There can also be a two-part track representing the party keeping watch for danger; it is a progressive track and each tick increases the encounter roll (which is part of the cost).
  \item There are generally two complication tracks: exhaustion and danger.
  \begin{itemize}
    \item The exhaustion track is a generally a six-part track and is progressive: when between 0 and 2 ticks have been made there is no effect, but when 3 or more ticks have been made the PCs' have a $-1$ penalty to all rolls, when 4 or more ticks have been made they only gain half of their Recovery when they heal, when 5 or more ticks have been made they suffer another $-1$ penalty to all rolls, and when all 6 ticks have been made they cannot heal at all.
    \item The danger track represents natural dangers and is a two-part track, when ticked it decreases the encounter rolls that are part of the cost.
    \item If the party is being pursued or followed there may be an additional track that represents being caught, the length is determined by how far behind their pursuers are and when it is completed the party is caught and has to deal with their pursuers.
  \end{itemize}
  \item The progress and exhaustion tracks are linked, when the progress track is ticked the exhaustion track ticks as well.
  \item For maintenance, the progress track ticks once each round and the track for keeping watch for danger resets every round.
  \item Finally, each round the cost is an encounter roll, based on the party's location the GM may roll for an encounter that the party has during the day.
    These should be a mix of combat and non-combat encounters, though in dangerous areas may be entirely combat encounters.
\end{enumerate}
